\documentclass[11pt,a4paper]{report}

\usepackage[utf8]{inputenc}
\usepackage[T1]{fontenc}
\usepackage[french]{babel}
\usepackage[margin=2.5cm]{geometry}
\usepackage{graphicx}
\usepackage{booktabs}
\usepackage{listings}
\usepackage{hyperref}

\begin{document}

% ====================================================================
\begin{titlepage}
  \centering
  Université de Montpellier --- Faculté des Sciences\\
  Master Informatique\\[2cm]

  {\Large HAI704I --- Architectures logicielles distribuées}\\[1cm]
  {\huge\bfseries TP 3 --- Services web REST}\\[2cm]

  Enseignant responsable : Bachar Rima\\
  Année universitaire 2026--2027\\
  \today\\[2cm]

  \begin{tabular}{llll}
    \toprule
    Nom & Prénom & N° étudiant & Parcours \\
    \midrule
    BAUDRAND & & & \\
    HAJEE & Neha & & \\
    \bottomrule
  \end{tabular}
\end{titlepage}

\tableofcontents

% ====================================================================
\chapter{Introduction}
\section{Rappel des objectifs du TP}
\section{Topologie de l'application}
\section{Dépôt Git et vidéo}

% ====================================================================
\chapter{A0 --- Consommer une API publique}
\section{Géoplateforme de l'IGN}
\section{Observer le protocole (GitHub, httpbingo)}
\section{Le même échange dans Postman}

% ====================================================================
\chapter{A1 --- Conception : du besoin aux ressources}
\section{Ressources identifiées}
\section{URI des ressources et des collections}
\section{Méthodes, sûreté, idempotence et codes de statut}
\section{Représentations}
\section{Relations entre ressources}
\section{Ce qui n'est pas une ressource}
\section{L'anti-modèle RPC argumenté}
\section{Ce que l'implémentation a fait changer}

% ====================================================================
\chapter{A2 --- L'API de lecture de l'hôtel}
\section{Recherche et consultation des offres}
\section{Enveloppe de collection}
\section{Séparation représentation / domaine / entité}
\section{La photo de la chambre}
\section{Validation}

% ====================================================================
\chapter{A3 --- Sémantique HTTP des erreurs}
\section{Centralisation en problem+json}
\section{Validation : 400 contre 422}
\section{Justification de chaque code de statut}
\section{L'en-tête Allow sur le 405}
\section{Validation}

% ====================================================================
\chapter{A4 --- L'agence : un client REST explicite}
\section{Projet et structure de l'agence}
\section{RestClient et lecture des statuts et en-têtes}
\section{Délais d'attente}
\section{Traduction des échecs distants}
\section{Séparation du transport et du métier}
\section{Validation}

% ====================================================================
\chapter{A5 --- Réserver}
\section{Création d'une réservation : 201 et Location}
\section{Authentification de l'agence}
\section{Offre déjà prise contre offre inexistante}
\section{Idempotence : le rejeu sans doublon}
\section{Concurrence : ETag et If-Match}
\section{Annulation}
\section{Validation}

% ====================================================================
\chapter{A6 --- Plusieurs hôtels et défaillance partielle}
\section{Liste d'hôtels configurable}
\section{Hôtel normal, lent et arrêté}
\section{Résultat partiel et incidents nommés}
\section{Refus métier contre défaillance}
\section{Code rendu quand l'amont tombe}
\section{Validation}

% ====================================================================
\chapter{A7 --- L'API de l'agence et le comparateur}
\section{L'agence comme serveur}
\section{Application de la marge}
\section{Identifiants d'offre uniques}
\section{Le comparateur : agrégation et classement}
\section{Deux agences, deux prix}
\section{Survie à une agence absente}
\section{Validation}

% ====================================================================
\chapter{A8 --- Application utilisable et persistance}
\section{De la mémoire à JPA/H2}
\section{Coût du changement et place de la frontière}
\section{Persistance après redémarrage}
\section{L'interface en ligne de commande}
\section{Tests automatisés}
\section{Chaîne de bout en bout}

% ====================================================================
\chapter{Parties bonus}
\section{B --- Base de données externe}
\section{C --- OpenAPI}
\section{D --- Cache et requêtes conditionnelles}
\section{E --- Hypermédia}
\section{F --- Interface graphique}

% ====================================================================
\chapter{Schémas}
\section{Diagramme de déploiement (processus et ports)}
\section{Diagramme de classes}
\section{Diagramme de séquence d'une réservation}

% ====================================================================
\chapter{Relevés HTTP commentés}
\section{Relevés de l'étape A0}
\section{Relevés de l'étape A3}
\section{Relevés de l'étape A5}

% ====================================================================
\chapter{Synthèse}
\section{Ce qui fonctionne}
\section{Ce qui ne fonctionne pas}
\section{Difficultés rencontrées}

% ====================================================================
\chapter{Regard critique}

% ====================================================================
\chapter{Usage des outils d'intelligence artificielle générative}
\section{Outils utilisés}
\section{Tâches concernées}
\section{Influence réelle sur le rendu}
\section{Vérifications et corrections}
\section{Ce qui est entièrement de nous}

% ====================================================================
\chapter*{Références}
\addcontentsline{toc}{chapter}{Références}

\end{document}
